\documentclass[12pt]{article}
\usepackage{amsmath, amssymb, enumitem, fancyhdr, graphicx, libertine, nth,
pgfplots, tikz, units, xcolor}
\pgfplotsset{compat=1.11}
\usepackage[margin=1in]{geometry}
\usepackage[libertine]{newtxmath}
\pagestyle{fancy}
\definecolor{solution-color}{HTML}{000080}
\chead{\emph{EC201 -- Problem Set 2}}

% Define new topic command
\newcommand{\topic}[1]{\medskip\begin{center}\large\textsc{#1}\normalsize
  \end{center}}

% Define question counter:
\newcounter{question}[section]
\newcommand{\question}[1]{\refstepcounter{question}
  \subsubsection*{Question~\thequestion~-- #1}}

% Subquestions are enumerate with letters
\newenvironment{subquestions}{
  \begin{enumerate}[label=(\roman*)]}{\end{enumerate}}

\begin{document}
\topic{Choice and Demand}
\question{Neutral Goods}
Revisiting the preferences from Question 6 on the previous problem set, suppose
that you enjoy good 1 but you are neutral about good 2 (your indifference
curves are vertical lines).
\begin{subquestions}
  \item Find the demand functions $x_1\left( p_1,p_2,m \right)$ and $x_2\left(
    p_1,p_2,m \right)$.
  \item Write down a utility function $u\left( x_1,x_2 \right)$ that would be
    consistent with these preferences.
\end{subquestions}

\question{Cobb-Douglas Demand}
Suppose you have the following utility function for goods 1 and 2:
\[
  u\left( x_1,x_2 \right) = 2 x_1^{\frac{1}{4}}x_2^{\frac{3}{4}}
\]
\begin{subquestions}
  \item Find the demand functions $x_1\left( p_1,p_2,m \right)$ and $x_2\left(
    p_1,p_2,m \right)$.
  \item The price of good 1 is \$1 and the price of good two is \$3. You have
    \$60 to spend. How much will you consume of goods 1 and 2?
  \item How much of goods 1 and 2 will you buy if you had \$120 and the prices
    stayed the same? Are the goods normal or inferior?
  \item If the price of good 2 doubled (to \$6), what would happen to your
    demand for good 1?
  \item Which of the following statements is true and why?
    \begin{itemize}
      \item Good 1 is a substitute to good 2.
      \item Good 1 is a complement to good 2.
      \item Good 1 is neither a substitute nor a complement to good 2.
    \end{itemize}
\end{subquestions}
\question{Income and Substitution Effects}
Suppose your utility function is $u\left( x_1,x_2 \right)=x_1^{\frac{2}{3}}x_2^{\frac{1}{3}}$. It can be shown that the resulting demand functions are $x_1\left( p_1,p_2,m \right)=\frac{2m}{3p_1}$ and $x_2\left( p_1,p_2,m \right)=\frac{m}{3p_2}$ (you should check this). Throughout this question, you have income of $m=30$.
\begin{subquestions}
\item If $p_1=1$ and $p_2=1$, how much of goods 1 and 2 would you buy? Sketch the budget line and optimal bundle on a graph.
\item If the price of good 1 increases to $p_1=2$ and the price of good 2 remains at $p_2$, how much of goods 1 and 2 would you buy? Sketch the new budget line and optimal bundle on your graph.
  \item Suppose that after the price of good 1 increased, the government gave you just enough money so that you could afford the original bundle, i.e.~the bundle in part(i). How much money would the government need to give you?
  \item If the government gave you this reimbursement, how much of goods 1 and 2 would you buy? Is this different from part (i)? Sketch your new budget line after the government handout and the optimal bundle.
  \item What is the income effect and substitution effect of the price increase of good 1 on the demand for good 1?
\end{subquestions}

\topic{Intertemporal Choice}
\question{Saving for Retirement}
There are two time periods. In period 1 you work and earn income $m_1$. In
period 2 you are retired and earn no income ($m_2=0$). The interest rate is $r$
and there is no inflation. Your consumption in time periods 1 and 2 are $c_1$
and $c_2$ respectively.

The lifetime budget constraint is therefore:
\[
  c_1 + \frac{c_2}{1+r} = m_1
\]
Your lifetime utility function is:
\[
  u\left( c_1,c_2 \right)=\frac{1}{2}c_1^2 + \frac{\beta}{2}c_2^2
\]
where $\beta$ is a constant.
\begin{subquestions}
  \item Sketch the budget constraint on a graph. Label the axis intercepts and
    the endowment.
  \item The marginal rate of substitution of consumption between the two
    periods is:
    \[
      MRS = - \frac{\frac{\partial u\left( c_1,c_2 \right)}{\partial c_1}}
                   {\frac{\partial u\left( c_1,c_2 \right)}{\partial c_2}}
    \]
    Calculate the $MRS$ for the lifetime utility function $u\left( c_1,c_2
    \right)$ provided above.
  \item Now we want to find exactly how much you will want to save in the first
    period. In what follows, \underline{assume that $\beta = \frac{1}{1+r}$}
    (this will make the algebra easier).

    At the optimum you will set the $MRS$ equal to the slope of the budget
    constraint, which is $-\left( 1+r \right)$.
    Therefore we have two equations to find $c_1$ and $c_2$:
    \begin{align}
      \frac{\frac{\partial u\left( c_1,c_2 \right)}{\partial c_1}}
           {\frac{\partial u\left( c_1,c_2 \right)}{\partial c_2}} &=
           1+r  \\
             c_1 + \frac{c_2}{1+r} &= m_1
    \end{align}
    Using these two equations, solve for $c_1$ and $c_2$ as a function of $m_1$
    and $r$.
  \item If the interest rate increases, will you consume more or less in the
    first period? To answer this you can assume that $m_1=1$ and $r$ ranges from
    0 to 1.
  \end{subquestions}
\newpage
\topic{Uncertainty}
% \question{Preferences for Risk}
% For each of the utility functions for wealth below, sketch them and state
% whether the consumer is risk averse, risk loving, or risk neutral.
% \begin{subquestions}
%   \item $u\left( W \right)=aW$, where $a>0$.
%   \item $u\left( W \right)=W^2$
%   \item $u\left( W \right)=\sqrt{W}$
% \end{subquestions}
\question{Insurance}
You have \$200,000 worth of valuable items in your home. There is a 5\% chance
that you will be burgled. Fortunately, you keep \$100,000 worth of your
valuables in a hidden safe that the burgulars won't find. The local insurance
company will insure your valuables for \$7,000. That is, in the event of a
burgulary the insurance company will give you \$100,000 (the amount the
burgulars stole).

Suppose your utility for wealth is $U\left( W \right)=\sqrt{W}$. So if you have
\$200,000 you utility is $\sqrt{200,000}$.

\begin{subquestions}
  \item What is your expected utility if you don't purchase insurance?
  \item What is your expected utility if you do purhcase insurance?
  \item Would you purchase insurance?
  \item How much would the insurance policy be if it were actuarially fair?
\end{subquestions}
\topic{Technology}
\question{The Technical Rate of Substitutition (TRS)}
Find the techinical rate of substitution for the following production function
(show all your work, including the partial derivatives to get the marginal
products):
\[
  f\left( x_1,x_2 \right)=x_1^{\frac{1}{3}}x_2^{\frac{2}{3}}
\]
Interpret what the $TRS$ means when $x_1 = 2$ and $x_2 = 4$.
\question{Returns to scale}
Show whether the following production functions exhibit constant returns to
scale, increasing returns to scale or decreasing returns to scale:
\begin{subquestions}
  \item $f\left( x_1,x_2 \right)=x_1x_2$.
  \item $f\left( x_1,x_2 \right)=x_1 + x_2$.
  \item $f\left( x_1,x_2 \right)=\min\left\{ x_1, x_2 \right\}$.
\end{subquestions}
\topic{Profit Maximization}
\question{Short-run profit maximization}
Suppose the production function is $y=f\left( x_1,x_2 \right)=x_1^{\frac{1}{2}}x_2^{\frac{1}{2}}$, the price of output is $p=2$ and the prices of each input/factor are $w_1=1$ and $w_2=1$. The firm is competitive and therefore has no control over input and output prices. Factor 2 is fixed in the short run at $x_2=1$. Find the amount of factor 1 that will maximimize the firm's profits. What profit does the firm make?
\end{document}
